\begin{conf-abstract}
{Origin of fluid emissions along crustal to lithospheric faults in an extensional setting, Betic Cordillera, Spain}
{Bérénice Cateland, Anne Battani, Matthias Brennwald, Benjamin Lefeuvre, Nicolas E. Beaudoin, Frédéric Mouthereau, Rolando Burga, Jose Benavente, Magali Pujol}
{Université de Pau et des pays de l'Adour}
{The Betic Cordillera (SE Spain) has experienced a complex geodynamic history involving crustal thinning related to slab retreat. This deep lithospheric dynamic was accompanied by alkaline to calc-alkaline volcanism and the exhumation of metamorphic domes. These processes feed an active fluid system, the respective contributions of which remain poorly constrained. However, understanding and quantifying these deep fluid circulations has major implications for geothermal exploration and natural gas migration. Several crustal to lithospheric-scale faults in the Internal Betics have accommodated most of the deformation. These active structures (Lorca earthquake, Mw 5.2) are associated with significant hydrothermalism, as indicated by water temperatures reaching up to 65\,°C, and may act as preferential pathways for fluid and heat migration. We sampled fluids (water and gas when bubbles were present) from 14 thermal springs with temperatures ranging from 15 to 65\,°C. Gas species were measured in the field using a portable mass spectrometer (MiniRUEDI) along these strike-slip faults. The Cádiz-Alicante Fault releases fluids rich in \ce{N2} (86–96\%), with \ce{CO2} contents between 1 and 5\%, and \ce{O2} between 0.5 and 10\%. Along the Alpujarras Fault, \ce{CO2} dominates (63–91\%) with lower amounts of \ce{N2} (5–28\%). The lithospheric Carboneras Fault emits fluids rich in \ce{N2} (91\%), with low proportions of \ce{O2} (3\%) and CO2 (3\%). Alhama de Murcia, the most active fault of the system, mainly releases \ce{CO2} (52\%) and \ce{N2} (43\%). This result helps us in discussing fluid origin in this extensive context.}
\end{conf-abstract}
